\section{Introduction}
\label{section:intro}

Large reasoning models (LRMs) \citep{jaech2024openai, rastogi2025magistral} increasingly exploit test-time compute by generating long chain-of-thought (CoT) traces. Challenging prompts are decoded over hundreds of thousands of tokens. A notable line of work, beginning with S1 \citep{muennighoff2025s1} and reinforced in subsequent papers \citep{ringel2025learning, jurayj2025your}, shows that appending \textit{wait} to the generation to increase test-time compute can improve reasoning performance. However, it is unclear whether such long reasoning traces are desired. Long reasoning traces not only significantly increase the resources for those hosting LRMs but also reduce user experience due to latency, especially for questions that intuitively do not require long reasoning. Moreover, recent studies \citep{wu2025more, hassid2025don, ghosal2025does, marjanovic2025deepseek} report that shorter thoughts are better, and continuing to append `wait' can induce oscillatory performance. Furthermore, it remains unclear whether different LRMs exhibit similar reasoning behaviors. 

These conflicting findings motivate a systematic re-examination of how lexical and structural properties of reasoning traces relate to reasoning performance. In this work, we evaluate the effectiveness of reasoning traces along multiple dimensions and uncover consistent patterns across LRMs. We analyze ten reasoning models with accessible reasoning traces on tasks spanning math and scientific reasoning (HARP, \citet{yue2024harp}, and GPQA-Diamond \citep{rein2024gpqa}), with the aim of providing systematic insight into what characterizes effective reasoning. 

\yunzhen{We begin by examining two properties that recent work suggests may drive reasoning performance: CoT length and review behaviors.} In the S1 approach, inserting \textit{wait} increases generation \Length and encourages \Review behaviors, including checking, verifying, or backtracking prior steps. These \Review behaviors are shown to be important to reasoning \citep{gandhi2025cognitive, chen2024not}. Therefore, we first investigate how \Length and \Review behaviors lead to reasoning improvement observed in \citet{muennighoff2025s1}. We define \rr as the fraction of \Review tokens within a CoT to isolate the effect of \Review from \Length. Using a conditional correlation analysis to isolate the question-level confounding factors, we find consistent patterns across models and datasets. \textit{Within the same question}, shorter reasoning traces are associated with higher accuracy, and lower \rr are associated with higher accuracy. 


We further hypothesize that \Length and \rr are surface proxies for underlying structural properties of the reasoning \citep{jiang2025makes} and we test one possible cause: the prevalence of failed reasoning branches. We therefore extract a \emph{reasoning graph} for each CoT. This representation allows for the evaluation of graph-level metrics. In particular, we focus on the \failedstep(\fsf): the fraction of steps belonging to failed exploratory branches.

Among graph-level features, \fsf emerges as a stronger and more stable predictor of correctness than CoT \Length or \rr, with consistent, significant correlations across difficulty strata and across all ten models on both math and scientific reasoning. These findings support measuring reasoning quality via the reasoning graph. Figure~\ref{fig:data_visualization} illustrates our annotation and the corresponding extracted reasoning graph.

Finally, we design two experiments to test causality. We first run a test-time intervention on AIME-25 and GPQA-Diamond: for each problem we sample 64 generations, rerank by each metric, and evaluate top-1 (pass@1) performance. \fsf-based selection yields the largest and most consistent gains, with up to 10\% accuracy improvement on AIME, while selection by \Length or \rr gives smaller benefits. Second, we intervene on the CoT directly via controlled editing: removing the failed branch substantially increases accuracy on incorrect traces. Together, these results provide causal evidence that \fsf is a strong lever for accuracy, that long failed branches bias subsequent exploration, and that current models do not fully “unsee” earlier mistakes when backtracking.


Our contributions can be summarized as:
\begin{itemize}
\item We conduct a wide-ranging conditional correlation test and show that, \emph{within the same question}, longer CoTs and higher \rr are negatively associated with accuracy. We measure a stronger correlation for harder questions.
\item We introduce a new reasoning-graph extraction method and define the \failedstep; \fsf robustly predicts correctness across models and difficulty strata, outperforming length and review ratio.
\item We design a test-time selection intervention providing causal evidence: \fsf-based reranking consistently outperforms baseline, \Length-, and \rr–based selection on AIME-25 and GPQA-Diamond.
\item We directly intervene in CoTs as a causal probe, revealing that failed attempts bias subsequent reasoning; \emph{removing} failed branches substantially improves accuracy.
\end{itemize}


% Our study begins by disentangling whether longer reasoning traces themselves lead to higher accuracy or whether the observed gains stem instead from review and verification behaviors embedded in the traces. Prior findings caution that excessive verification can amount to overthinking \citep{chen2024not, su2025between}, raising the possibility that more review may hinder rather than help.
